\documentclass[11pt,a4paper]{article}

\usepackage{fontspec}
\usepackage[french]{babel}
\usepackage{amsmath,amssymb,amsthm,mathtools}
\usepackage{geometry}
\usepackage{xcolor}
\usepackage{fancyhdr}
\usepackage{enumitem}
\usepackage{hyperref}
\usepackage{titlesec}
\usepackage{microtype}

\geometry{margin=2.2cm}
\setmainfont{DejaVu Serif}
\setsansfont{DejaVu Sans}
\setmonofont{DejaVu Sans Mono}

\definecolor{accent}{HTML}{1F4E79}
\definecolor{lightaccent}{HTML}{EAF2F8}

\pagestyle{fancy}
\fancyhf{}
\lhead{\small Cours Deep Reinforcement Learning}
\chead{\small Classe : 4IA}
\rhead{\small Année : 2026}
\lfoot{\small Enseignant : Mourad Zéraï, PhD - ESPRIT School of Engineering}
\cfoot{}
\rfoot{\thepage}
\setlength{\headheight}{14pt}

\titleformat{\section}{\large\bfseries\color{accent}}{\thesection.}{0.5em}{}
\titleformat{\subsection}{\normalsize\bfseries\color{accent}}{\thesubsection.}{0.5em}{}

\newtheorem*{definition}{Définition}
\newtheorem*{rappel}{Rappel}
\newtheorem*{hypothese}{Hypothèse}
\newtheorem*{propriete}{Propriété}

\newcommand{\E}{\mathbb{E}}
\newcommand{\Pbb}{\mathbb{P}}
\newcommand{\R}{\mathbb{R}}
\newcommand{\given}{\,|\,}

\begin{document}

\begin{center}
    {\LARGE\bfseries Dérivation de l'équation de Bellman pour la fonction état-action $q_{\pi}$}\\[0.4em]
    {\large Note de cours rigoureuse et détaillée}
\end{center}

\vspace{0.6em}

\section{Objectif}

Le but de cette note est de dériver, pas à pas, l'équation de Bellman vérifiée par la fonction état-action associée à une politique $\pi$, à savoir
\[
q_{\pi}(s,a) = \E_{\pi}\!\left[R_{t+1} + \gamma q_{\pi}(S_{t+1},A_{t+1}) \middle\vert S_t=s, A_t=a\right].
\]
Nous justifierons chaque passage à partir des définitions fondamentales, de la décomposition du retour, de la linéarité de l'espérance conditionnelle, de la loi de l'espérance totale et de la propriété de Markov.

\section{Cadre probabiliste et notations}

\begin{hypothese}
On considère un processus de décision markovien (MDP) à temps discret. Pour chaque date $t \in \mathbb{N}$ :
\begin{itemize}[leftmargin=1.6em]
    \item $S_t$ désigne l'état à l'instant $t$ ;
    \item $A_t$ désigne l'action choisie à l'instant $t$ ;
    \item $R_{t+1}$ désigne la récompense reçue après avoir exécuté $A_t$ en $S_t$ ;
    \item $\gamma \in [0,1)$ est le facteur d'actualisation.
\end{itemize}
La politique $\pi$ est une loi conditionnelle sur les actions :
\[
\pi(a\given s) = \Pbb(A_t=a \given S_t=s).
\]
\end{hypothese}

\begin{definition}[Retour]
Le retour à partir de l'instant $t$ est la variable aléatoire
\[
G_t = \sum_{k=0}^{\infty} \gamma^k R_{t+k+1},
\]
sous réserve que cette série soit bien définie, ce qui est le cas par exemple si les récompenses sont bornées et si $\gamma \in [0,1)$.
\end{definition}

\begin{definition}[Fonction état-action]
La fonction état-action sous la politique $\pi$ est définie par
\[
q_{\pi}(s,a) = \E_{\pi}[G_t \given S_t=s, A_t=a].
\]
Autrement dit, $q_{\pi}(s,a)$ est le retour espéré quand on se trouve dans l'état $s$, que l'on impose l'action initiale $a$, puis que l'on suit ensuite la politique $\pi$.
\end{definition}

\section{Formule de départ}

Partons de la définition :
\[
q_{\pi}(s,a) = \E_{\pi}[G_t \given S_t=s, A_t=a].
\]
La première étape consiste à décomposer le retour $G_t$.

\subsection{Décomposition récursive du retour}

\begin{propriete}
Pour tout $t$,
\[
G_t = R_{t+1} + \gamma G_{t+1}.
\]
\end{propriete}

\paragraph{Justification.}
En partant de la définition du retour,
\[
G_t = R_{t+1} + \gamma R_{t+2} + \gamma^2 R_{t+3} + \cdots.
\]
On factorise tous les termes à partir du second par $\gamma$ :
\[
G_t = R_{t+1} + \gamma\left(R_{t+2} + \gamma R_{t+3} + \gamma^2 R_{t+4} + \cdots\right).
\]
Or, par définition même du retour à l'instant suivant,
\[
G_{t+1} = R_{t+2} + \gamma R_{t+3} + \gamma^2 R_{t+4} + \cdots.
\]
Donc
\[
G_t = R_{t+1} + \gamma G_{t+1}.
\]

\section{Substitution dans la définition de $q_{\pi}$}

En remplaçant $G_t$ par son expression récursive dans la définition de $q_{\pi}$, on obtient
\begin{align*}
q_{\pi}(s,a)
&= \E_{\pi}[G_t \given S_t=s, A_t=a] \\
&= \E_{\pi}[R_{t+1} + \gamma G_{t+1} \given S_t=s, A_t=a].
\end{align*}

\subsection{Utilisation de la linéarité de l'espérance conditionnelle}

\begin{propriete}[Linéarité]
Si $X$ et $Y$ sont intégrables et si $c \in \R$, alors
\[
\E[X + cY \given \mathcal{F}] = \E[X \given \mathcal{F}] + c\,\E[Y \given \mathcal{F}].
\]
\end{propriete}

On applique cette propriété avec
\[
X = R_{t+1}, \qquad Y = G_{t+1}, \qquad c = \gamma,
\]
et avec la tribu engendrée par l'information ``$S_t=s, A_t=a$''. Il vient
\begin{align*}
q_{\pi}(s,a)
&= \E_{\pi}[R_{t+1} \given S_t=s, A_t=a]
+ \gamma\,\E_{\pi}[G_{t+1} \given S_t=s, A_t=a].
\end{align*}

Cette forme est correcte, mais elle ne fait pas encore apparaître explicitement l'état suivant $S_{t+1}$ ni l'action suivante $A_{t+1}$. Il faut donc raffiner la seconde espérance conditionnelle.

\section{Conditionnement supplémentaire sur la transition suivante}

On s'intéresse au terme
\[
\E_{\pi}[G_{t+1} \given S_t=s, A_t=a].
\]
Pour faire apparaître l'état et l'action au temps $t+1$, on utilise la loi de l'espérance totale sous forme conditionnelle.

\begin{propriete}[Loi de l'espérance totale]
Si $X$ est intégrable et si $Y,Z$ sont des variables aléatoires, alors
\[
\E[X \given Y] = \E\big[\E[X \given Y,Z] \given Y\big].
\]
Dans un cadre discret, cela s'écrit aussi sous forme de somme pondérée par les probabilités conditionnelles.
\end{propriete}

En appliquant cette formule avec
\[
X = G_{t+1}, \qquad Y=(S_t,A_t), \qquad Z=(S_{t+1},A_{t+1}),
\]
on obtient
\[
\E_{\pi}[G_{t+1} \given S_t=s, A_t=a]
=
\E_{\pi}\!\left[\E_{\pi}[G_{t+1} \given S_t=s, A_t=a, S_{t+1}, A_{t+1}] \middle\vert S_t=s, A_t=a\right].
\]

Dans un espace d'états et d'actions discret, cette identité devient
\begin{align*}
\E_{\pi}[G_{t+1} \given S_t=s, A_t=a]
&= \sum_{s',a'} \Pbb_{\pi}(S_{t+1}=s', A_{t+1}=a' \given S_t=s, A_t=a) \\
&\hspace{2em} \times \E_{\pi}[G_{t+1} \given S_t=s, A_t=a, S_{t+1}=s', A_{t+1}=a'].
\end{align*}

\section{Utilisation de la propriété de Markov}

Le point clé est maintenant le suivant.

\begin{propriete}[Propriété de Markov]
Dans un MDP, conditionnellement à l'état courant et à l'action courante, la loi de la transition suivante ne dépend pas du passé complet. Plus précisément, la loi de $(S_{t+1},R_{t+1})$ conditionnée par l'histoire ne dépend que de $(S_t,A_t)$.

De plus, à partir de l'instant $t+1$, si l'on connaît $S_{t+1}$ et l'action suivante $A_{t+1}$, l'espérance du retour futur ne dépend plus de $(S_t,A_t)$.
\end{propriete}

Ainsi,
\[
\E_{\pi}[G_{t+1} \given S_t=s, A_t=a, S_{t+1}=s', A_{t+1}=a']
=
\E_{\pi}[G_{t+1} \given S_{t+1}=s', A_{t+1}=a'].
\]
Par définition de la fonction état-action, mais appliquée à l'instant $t+1$ au lieu de $t$,
\[
\E_{\pi}[G_{t+1} \given S_{t+1}=s', A_{t+1}=a'] = q_{\pi}(s',a').
\]
Donc
\[
\E_{\pi}[G_{t+1} \given S_t=s, A_t=a]
= \sum_{s',a'} \Pbb_{\pi}(s',a' \given s,a)\, q_{\pi}(s',a').
\]
Ici, pour alléger l'écriture,
\[
\Pbb_{\pi}(s',a' \given s,a)
:= \Pbb_{\pi}(S_{t+1}=s', A_{t+1}=a' \given S_t=s, A_t=a).
\]

\section{Retour à l'équation principale}

En remplaçant ce résultat dans l'expression précédente de $q_{\pi}(s,a)$, on trouve
\begin{align*}
q_{\pi}(s,a)
&= \E_{\pi}[R_{t+1} \given S_t=s, A_t=a]
+ \gamma \sum_{s',a'} \Pbb_{\pi}(s',a' \given s,a)\, q_{\pi}(s',a').
\end{align*}

Cette écriture est déjà une forme correcte de l'équation de Bellman pour $q_{\pi}$. Il reste à expliciter la probabilité conjointe sous la politique $\pi$.

\section{Factorisation de la loi conjointe sous la politique}

Sous une politique stationnaire $\pi$, une fois l'état suivant $S_{t+1}=s'$ réalisé, l'action suivante $A_{t+1}$ est tirée selon la loi $\pi(\cdot \given s')$. Par conséquent,
\[
\Pbb_{\pi}(S_{t+1}=s', A_{t+1}=a' \given S_t=s, A_t=a)
=
\Pbb(S_{t+1}=s' \given S_t=s, A_t=a)\,\pi(a'\given s').
\]

Si l'on note
\[
p(s'\given s,a) = \Pbb(S_{t+1}=s' \given S_t=s, A_t=a),
\]
alors
\[
\Pbb_{\pi}(s',a' \given s,a) = p(s'\given s,a)\,\pi(a'\given s').
\]
D'où
\[
q_{\pi}(s,a)
= \E_{\pi}[R_{t+1} \given S_t=s, A_t=a]
+ \gamma \sum_{s',a'} p(s'\given s,a)\,\pi(a'\given s')\, q_{\pi}(s',a').
\]

\section{Incorporation explicite de la récompense immédiate}

Dans un MDP général, la récompense peut dépendre de $s$, $a$, $s'$ et éventuellement de sa propre loi conditionnelle. Pour écrire l'équation la plus standard, on introduit la dynamique jointe
\[
p(s',r \given s,a) := \Pbb(S_{t+1}=s', R_{t+1}=r \given S_t=s, A_t=a).
\]
Alors, par définition de l'espérance d'une variable discrète,
\[
\E_{\pi}[R_{t+1} \given S_t=s, A_t=a]
= \sum_{s',r} p(s',r \given s,a)\,r.
\]
En remplaçant ce terme, on obtient
\[
q_{\pi}(s,a)
= \sum_{s',r} p(s',r \given s,a)\,r
+ \gamma \sum_{s',a'} p(s'\given s,a)\,\pi(a'\given s')\, q_{\pi}(s',a').
\]

On peut fusionner les deux sommes en une seule écriture plus compacte :
\[
q_{\pi}(s,a)
= \sum_{s',r} p(s',r \given s,a)
\left[r + \gamma \sum_{a'} \pi(a'\given s') q_{\pi}(s',a')\right].
\]
C'est la forme classique de l'équation de Bellman pour la fonction état-action.

\section{Forme compacte via la fonction de valeur d'état}

Rappelons que la fonction de valeur d'état est définie par
\[
v_{\pi}(s') = \E_{\pi}[G_{t+1} \given S_{t+1}=s'] = \sum_{a'} \pi(a'\given s') q_{\pi}(s',a').
\]
Cette identité provient encore d'une loi de l'espérance totale conditionnée sur l'action choisie dans l'état $s'$.

On peut alors écrire l'équation de Bellman pour $q_{\pi}$ sous la forme
\[
q_{\pi}(s,a)
= \sum_{s',r} p(s',r \given s,a)\left[r + \gamma v_{\pi}(s')\right].
\]
Cette écriture est souvent la plus lisible lorsque l'on veut faire apparaître le lien entre $q_{\pi}$ et $v_{\pi}$.

\section{Résumé logique de la dérivation}

La dérivation repose sur la chaîne d'idées suivante :
\begin{enumerate}[leftmargin=1.8em]
    \item on part de la définition
    \[
    q_{\pi}(s,a)=\E_{\pi}[G_t \given S_t=s,A_t=a];
    \]
    \item on utilise la décomposition récursive du retour
    \[
    G_t = R_{t+1} + \gamma G_{t+1};
    \]
    \item on applique la linéarité de l'espérance conditionnelle ;
    \item on introduit $(S_{t+1},A_{t+1})$ via la loi de l'espérance totale ;
    \item on remplace l'espérance conditionnelle du retour futur par $q_{\pi}(s',a')$ grâce à la propriété de Markov et à la définition de $q_{\pi}$ ;
    \item on développe la loi jointe sous la politique par
    \[
    \Pbb_{\pi}(s',a'\given s,a)=p(s'\given s,a)\pi(a'\given s');
    \]
    \item on obtient finalement l'équation de Bellman pour la fonction état-action.
\end{enumerate}

\section{Formule finale}

Les formes usuelles à retenir sont donc :
\[
q_{\pi}(s,a) = \E_{\pi}\!\left[R_{t+1} + \gamma q_{\pi}(S_{t+1},A_{t+1}) \middle\vert S_t=s, A_t=a\right],
\]
\[
q_{\pi}(s,a)
= \sum_{s',r} p(s',r \given s,a)
\left[r + \gamma \sum_{a'} \pi(a'\given s') q_{\pi}(s',a')\right],
\]
et encore
\[
q_{\pi}(s,a)
= \sum_{s',r} p(s',r \given s,a)\left[r + \gamma v_{\pi}(s')\right].
\]

\vspace{0.8em}

\noindent\textbf{Remarque pédagogique.} La différence essentielle avec la dérivation de $v_{\pi}$ est que l'on conditionne ici au départ sur \emph{un couple} $(S_t,A_t)$, ce qui fait apparaître naturellement, à l'étape suivante, le couple $(S_{t+1},A_{t+1})$ et donc une moyenne sur la politique seulement à partir du temps $t+1$.

\end{document}
